\section*{The Bangalore Astronomical Society}
\begin{wrapfigure}{l}{0.2\textwidth}
\vspace{-15pt}
\centering
\includegraphics[width=0.175\textwidth]{BASlogo.eps}
\vspace{-30pt}
\end{wrapfigure}
The makers of this compilation acknowledge the Bangalore Astronomical
Society (BAS) for the inspiration behind this idea. In particular, the
makers thank the council members of the BAS during 2013.\\
\vspace{10pt}
\section*{The Digitized Sky Survey}
The images used in this compilation come from the Digitized Sky Survey
plates, in particular, those from the POSS-II and UKSTU surveys.

The Digitized Sky Survey was produced at the Space Telescope Science
Institute under U.S. Government grant NAG W-2166. The images of these
surveys are based on photographic data obtained using the Oschin
Schmidt Telescope on Palomar Mountain and the UK Schmidt
Telescope. The plates were processed into the present compressed
digital form with the permission of these institutions.

The Second Palomar Observatory Sky Survey (POSS-II) was made by the
California Institute of Technology with funds from the National
Science Foundation, the National Aeronautics and Space Administration,
the National Geographic Society, the Sloan Foundation, the Samuel
Oschin Foundation, and the Eastman Kodak Corporation. The Oschin
Schmidt Telescope is operated by the California Institute of
Technology and Palomar Observatory.

The UK Schmidt Telescope was operated by the Royal Observatory
Edinburgh, with funding from the UK Science and Engineering Research
Council (later the UK Particle Physics and Astronomy Research
Council), until 1988 June, and thereafter by the Anglo-Australian
Observatory. The blue plates of the southern Sky Atlas and its
Equatorial Extension (together known as the SERC-J), the near-IR
plates (SERC-I), as well as the Equatorial Red (ER), and the Second
Epoch [red] Survey (SES) were all taken with the UK Schmidt telescope
at the AAO.

The images themselves were downloaded from the Mikulski Archive for
Space Telescopes (MAST; \url{http://archive.stsci.edu/}).

The makers thank the DSS for making sky imagery freely available for
non-profit activities, and also thank MAST for the excellent web
service provided by them.

\section*{Deep-Sky Object Data}

The makers thank Dr. Wolfgang Steinicke for providing the Revised NGC
/ IC catalog under terms making it free for non-commercial use.

The Dreyer and SAC descriptions, wherever available, come from the
Saguaro Astronomy Club database. The makers thank the Saguaro
Astronomy Club for providing their compilation for free non-commercial
use.
\clearpage
\section*{KStars and other open-source tools}
\begin{wrapfigure}{l}{0.2\textwidth}
\centering
\includegraphics[width=0.195\textwidth]{kstars.pdf}
\end{wrapfigure}
The makers particulary thank, the developers of KStars,
(\url{http://edu.kde.org/kstars}) the software that made the rendition
of star maps used in this compilation possible and made available, in
an easy form, the data used in this compilation. KStars was also used
to fetch appropriate DSS URLs for the objects. KStars is a
cross-platform astronomy software licensed under the GNU General
Public License v2 (\url{https://www.gnu.org/licenses/gpl-2.0}). It
qualifies as free software.

The typesetting of the charts was done using
\LaTeX. \texttt{xmlstarlet} was used to parse XML for object
descriptions generated by KStars. Inkscape and ImageMagick were used
to convert between graphics formats. Inkscape was also used to
generate several of the graphics used here. Several tools from the
standard GNU suite, such as BASH, \texttt{sed} and \texttt{awk} proved
very useful.

\textbf{This compilation was generated using only free and open source
software.}
